\section{Results}
\label{sec:results}

\subsection{Dataset and Testability}
\label{sec:res-testability}

No simulation failed. All healthy cases are compliant, and so are \SI{79.0}{\percent} of all
cases of \cA\ and \SI{69.8}{\percent} of \cB, although fewer than \SI{8}{\percent} are healthy:
most injected faults do not take the circuit out of specification
(Fig.~\ref{fig:compliance}). With one component \SI{5}{\percent} off nominal, 93 to
\SI{96}{\percent} of the circuits remain compliant; with \SI{50}{\percent}, 54 to
\SI{66}{\percent}. In \cA\ no fault of the DRL integrator or of its input buffer violates a
specification, and neither does a CMRR of the integrated amplifier degraded to
\SI{50}{\decibel}, because the common-mode rejection of the system is dominated by the DRL path.

\begin{figure}[!t]
\centering
\includegraphics[width=\columnwidth]{compliance_by_magnitude}
\caption{Circuits that still meet every specification when one passive component deviates from
nominal by the given percentage. Mean over all passive components, 200 cases each.}
\label{fig:compliance}
\end{figure}

The electrode changes what the self-test sees. The gain that a healthy \cA\ measures through
$V_\mathrm{cal}$ at \SI{10}{\hertz}, relative to nominal, lies between 0.90 and 1.02 with gel and
solid-metal electrodes, between 0.59 and 1.00 with the polymer, and between 0.05 and 0.97 with
the fabrics, whose megaohm-level impedance forms a divider with the input network. The circuit
is nonetheless compliant, since the standard tests the input impedance with
\SI{620}{\kilo\ohm}.

\begin{table}[!t]
\caption{Testability of the Main Feature Set (Univariate Criterion)}
\label{tab:testability}
\centering
\setlength{\tabcolsep}{3pt}
\begin{tabular}{@{}p{0.50\columnwidth}cccc@{}}
\toprule
 & \multicolumn{2}{c}{\cA} & \multicolumn{2}{c}{\cB} \\
\cmidrule(lr){2-3}\cmidrule(l){4-5}
Electrodes & All & No porous & All & No porous \\
\midrule
Non-compliant conditions$^{\mathrm{a}}$ indistinguishable from healthy & 9 & 1 & 16 & 6 \\
Non-compliant conditions$^{\mathrm{a}}$ missed by the limit test & 15 & 3 & 35 & 8 \\
Testability rank (\SI{10}{\percent} deviation) & 2 & 4 & 2 & 3 \\
Components located without ambiguity & 4/23 & 5/23 & 2/29 & 5/29 \\
Mean number of components each one can be mistaken for & 3.9 & 2.5 & 8.6 & 5.2 \\
\bottomrule
\multicolumn{5}{@{}p{0.97\columnwidth}@{}}{\footnotesize $^{\mathrm{a}}$Fault conditions in
which most cases are non-compliant, out of 293 (\cA) and 307 (\cB).}
\end{tabular}
\end{table}


Table~\ref{tab:testability} summarizes the analysis of Section~\ref{sec:testability}. First,
some non-compliant faults cannot be seen. In \cA, eight of the nine invisible conditions are gain
errors of 5 to \SI{20}{\percent}, confused with the attenuation of a high-impedance electrode;
without porous electrodes only the open output resistor of the DRL stage remains. In \cB\ six
remain, all in the DRL stage. Second, few components can be located without ambiguity, and the
groups are predictable. The sensitivities~\eqref{eq:sensitivity} alone give, for \cA,
R5+R6+R11+R12 (the resistors that set the gain of the two stages), C5+R10 (high-pass filter),
R13+R14 (low-pass filter) and R15+R16 (reference divider); for \cB, one group with the nine
resistors that set the gain and three groups of two. Third, the structure depends on the
architecture: \cA\ has fewer invisible faults and each component can be mistaken for fewer
others.

\subsection{Compliance Decision}
\label{sec:res-decision}

\begin{table}[!t]
\caption{Compliance Decision From the Main Feature Set. Mean of Three Repetitions and
\SI{95}{\percent} Confidence Interval, in Percent}
\label{tab:decision}
\centering
\setlength{\tabcolsep}{3.5pt}
\begin{tabular}{@{}llcc@{}}
\toprule
Circuit & Method & Escapes & False rejects \\
\midrule
A & Limit test & 22.5 (21.4--23.6) & 20.0 (19.8--20.3) \\
 & Specification regression & 8.4 (6.7--10.2) & 1.5 (1.3--1.6) \\
 & Classifier & 7.9 (6.4--9.4) & 0.7 (0.4--1.0) \\
 & Classifier, \SI{1}{\percent} escape target & 0.9 (0.8--1.0) & 7.8 (4.8--10.8) \\
\addlinespace
B & Limit test & 35.3 (34.7--35.8) & 17.0 (16.5--17.6) \\
 & Specification regression & 11.4 (10.0--12.8) & 2.4 (1.9--2.9) \\
 & Classifier & 11.2 (10.3--12.1) & 1.2 (1.1--1.2) \\
 & Classifier, \SI{1}{\percent} escape target & 1.2 (0.6--1.7) & 54.0 (48.0--60.0) \\
\addlinespace
\multicolumn{4}{@{}l}{\emph{Without porous dry electrodes}} \\
A & Classifier & 4.2 (3.5--4.9) & 0.7 (0.4--0.9) \\
 & Classifier, \SI{1}{\percent} escape target & 1.0 (0.9--1.1) & 1.9 (1.3--2.5) \\
B & Classifier & 6.7 (6.2--7.1) & 1.0 (0.6--1.4) \\
 & Classifier, \SI{1}{\percent} escape target & 1.0 (0.7--1.3) & 48.7 (39.4--58.1) \\
\bottomrule
\end{tabular}
\end{table}


\begin{figure}[!t]
\centering
\includegraphics[width=\columnwidth]{operating_points}
\caption{Operating points of the four decision methods on the 26 features. Both axes are
logarithmic; confidence intervals are in Table~\ref{tab:decision}.}
\label{fig:operating}
\end{figure}

Table~\ref{tab:decision} and Fig.~\ref{fig:operating} give the result. The limit test accepts
between a fifth and a third of the non-compliant circuits and rejects a fifth of the compliant
ones: a wide healthy envelope, widened further by the electrodes, and a narrow compliance region
do not coincide. The specification regression and the classifier are close to each other at
their default operating points; the root-mean-square error of the predicted specifications is
between 8 and \SI{36}{\percent} of the limit, the largest for the input-impedance test.

An instrument needs a low escape rate. Lowering the threshold of the classifier brings escapes
on \cA\ to \SI{0.9}{\percent} at \SI{7.8}{\percent} of false rejects, and at only
\SI{1.9}{\percent} without porous electrodes. On \cB\ the same escape rate costs half of the
compliant circuits: the escapes that remain are the DRL faults that the measurements do not see,
as the testability analysis anticipated, and no threshold recovers a fault that leaves no trace.
A guard band on the regression is not a usable operating point (more than \SI{90}{\percent} of
false rejects at the same target). Telling the classifier which electrode type is connected does
not help (\SI{8.3}{\percent} of false rejects on \cA): the variation that matters is between
subjects, not between materials. No single measurement set suffices; with the escape target the
pulse response alone gives \SI{44}{\percent} of false rejects on \cA\ and every other set more
than \SI{75}{\percent}.

\subsection{Functional Against Percentage Severity}
\label{sec:res-severity}

A detector trained to recognize that a component is out of tolerance, scored against compliance,
rejects \SI{36.8}{\percent} of the compliant circuits of \cA\ and \SI{36.5}{\percent} of \cB,
with 7.9 and \SI{8.3}{\percent} of escapes. Trained on compliance, the same model rejects 1.6
and \SI{2.2}{\percent}, with 4.1 and \SI{7.6}{\percent} of escapes. The false rejects of the
percentage-trained detector are dominated by parametric faults that are real but harmless. (The
escapes of the functional detector are lower than in Table~\ref{tab:decision} because this
experiment weights the classes, which moves the operating point.)

\subsection{Location and Origin}
\label{sec:res-location}

There are \num{13364} non-compliant cases of \cA\ over 23 components and \num{20033} of \cB\ over
29. At the level of the a priori groups the best models reach a macro F1 of 0.87 on \cA\
(gradient boosting) and 0.76 on \cB\ (random forest), and the true group is among the three most
probable in more than \SI{99}{\percent} of the cases. At component level the macro F1 falls to
0.74 and 0.61, and the errors concentrate in the predicted groups: in \cA\ the recall is 36 to
\SI{41}{\percent} for the halves of the gain resistor and for the low-pass resistors, against 88
to \SI{100}{\percent} for most components outside a group. The recall of a component decreases
with the number of components that the testability analysis predicted it could be mistaken for
(Spearman $\rho=-0.42$, $p=0.045$ for \cA; $\rho=-0.51$, $p=0.005$ for \cB). A convolutional
network on the raw pulse response stays below the tabular models (0.56 and 0.32).

For origin, the macro F1 is 0.89 on \cA\ and 0.87 on \cB. A case with nothing to act on is
recognized in \SI{99}{\percent} of the cases and a non-compliant circuit in 94 and
\SI{92}{\percent}; circuit faults are almost never taken for electrode problems. Electrode
problems are recognized in only 64 and \SI{61}{\percent} of the cases. An open protection
resistor is always recognized and a detached electrode in 75 and \SI{89}{\percent} of the cases,
but a degraded contact only in half of the cases in \cA\ and \SI{41}{\percent} in \cB: a contact
impedance five or twenty times higher than normal for one subject falls within what healthy
contacts show on other subjects. Without porous electrodes the recall of the electrode class
rises to 76 and \SI{74}{\percent}.

\subsection{Robustness}
\label{sec:res-robustness}

\begin{figure}[!t]
\centering
\includegraphics[width=\columnwidth]{unseen_magnitudes}
\caption{Location of the cause for parametric faults whose magnitudes were held out of training.
The model names the component; the group bars count a prediction as correct when it falls in the
ambiguity group of the true component.}
\label{fig:unseen}
\end{figure}

The compliance decision generalizes to fault magnitudes absent from training: its balanced
accuracy falls by less than 2 points with respect to a model trained with every magnitude,
except for \SI{50}{\percent} deviations on \cA\ (6 points). Component-level location does not
(Fig.~\ref{fig:unseen}): correctly located components fall from between 21 and \SI{70}{\percent}
when the magnitude was in training to between 2 and \SI{33}{\percent}. Within a group of
components with parallel sensitivities the model can only tell the members apart by the size of
the effect, that is, by memorizing the magnitudes it was shown. Counted at group level, the same
predictions remain correct in 84 to \SI{99}{\percent} of the cases on \cA\ and 79 to
\SI{96}{\percent} on \cB.

When the models are trained at the noise level and resolution at which they are used, results
barely change from 0.1 to \SI{10}{\milli\volt} of noise and from 8 to 16 bits. Trained at
\SI{1}{\milli\volt} and used at \SI{10}{\milli\volt}, escapes on \cA\ rise from 4 to
\SI{10}{\percent} and false rejects on \cB\ from 2 to \SI{19}{\percent}. A truncated Gaussian
tolerance distribution changes nothing; \SI{2}{\percent} resistors and \SI{10}{\percent}
capacitors raise escapes from 4.1 to \SI{7.1}{\percent} on \cA\ and from 7.6 to
\SI{11.4}{\percent} on \cB, and leave group-level location unchanged on \cA\ (F1 0.87) but not on
\cB\ (0.45). Raising the common-mode tone from 0.1 to \SI{1}{\volt}, or making it ten times
longer, does not make the DRL faults detectable: about half of them escape whatever the setting.
The limitation is not one of signal-to-noise ratio.

\subsection{Minimal Set of Measurements}
\label{sec:res-minimal}

\begin{figure}[!t]
\centering
\includegraphics[width=\columnwidth]{cost_performance}
\caption{Balanced accuracy of the compliance decision on the test part as self-test measurements
are added one at a time by greedy selection on the validation part. One line per repetition.}
\label{fig:cost}
\end{figure}

The calibration pulse is always selected first (Fig.~\ref{fig:cost}). For the decision on \cA,
three measurements taking \SI{3}{\second} (the pulse, one lead-off tone and the differential tone
at \SI{150}{\hertz}) are within one point of the complete self-test of \SI{97}{\second}, in the
three repetitions. \cB\ needs a fourth, a dc reading of the DRL output, which is the measurement
that sees its DRL faults. For location on \cA, the pulse, the \SI{150}{\hertz} tone and the dc
reading of the instrumentation-amplifier output (\SI{2}{\second}) suffice in two of the three
repetitions; \cB\ needs five or six measurements. The \SI{0.05}{\hertz} tones, 80 of the
\SI{97}{\second}, are never among the first three selected.

Tables and figures for the experiments summarized in this section are in the supplementary
material.
